\documentclass[11pt]{article}
\usepackage[T1]{fontenc}
\usepackage[utf8]{inputenc}
\usepackage{amsmath,amssymb,amsthm,mathtools}
\usepackage{fullpage}
\usepackage{enumitem}
\usepackage{bbm}
\usepackage{hyperref}

% --- theorem environments
\newtheorem{theorem}{Theorem}[section]
\newtheorem{proposition}[theorem]{Proposition}
\theoremstyle{remark}
\newtheorem{remark}[theorem]{Remark}
\theoremstyle{definition}
\newtheorem{definition}[theorem]{Definition}
\newtheorem{lemma}[theorem]{Lemma}
\newtheorem{corollary}[theorem]{Corollary}
\newtheorem{conjecture}[theorem]{Conjecture}
\newtheorem{problem}[theorem]{Problem}

% --- macros
\newcommand{\C}{\mathbb{C}}
\newcommand{\D}{\mathbb{D}}
\newcommand{\ip}[2]{\left\langle #1,\,#2\right\rangle}
\newcommand{\norm}[1]{\left\|#1\right\|}
\newcommand{\opnorm}[1]{\left\|#1\right\|_2}
\newcommand{\Herm}{\mathrm{Herm}}
\newcommand{\diag}{\mathrm{diag}}
\newcommand{\spec}{\mathrm{spec}}
\newcommand{\tr}{\mathrm{tr}}
\newcommand{\Id}{\mathbbm{1}}
\newcommand{\W}{W}
\newcommand{\kappaopt}{\chi}
\newcommand{\adj}{^{\ast}}
\newcommand{\rank}{\operatorname{rank}}

\DeclareMathOperator{\cond}{cond}
\DeclareMathOperator{\supp}{supp}

\title{\textbf{Crouzeix's Conjecture for Nilpotent Dimension $3$}}
\author{Draft for circulation}
\date{\today}

\begin{document}
\maketitle

% ----------------------------------------------------------------
\section{Introduction: Crouzeix ratio and the nilpotent $3\times 3$ case}
\label{sec:intro}

For $T\in\C^{n\times n}$, recall the numerical range (field of values)
\[
\W(T):=\{\ip{Tx}{x}:\ x\in\C^n,\ \norm{x}=1\}.
\]
It is classical that $\W(T)$ is a compact convex subset of $\C$ containing $\spec(T)$.

A central theme in functional calculus is to compare $\opnorm{p(T)}$ to the sup norm of $p$
on a set in the complex plane. For the numerical range this leads to the following quantity.

\begin{definition}[Crouzeix ratio]
\label{def:crouzeix-ratio}
For $T\in\C^{n\times n}$ define
\[
\mathcal C(T)
\;:=\;
\sup_{p\not\equiv 0}\,
\frac{\opnorm{p(T)}}{\max_{z\in \W(T)}|p(z)|}.
\]
\end{definition}

Crouzeix's conjecture asserts that $\mathcal C(T)\le 2$ for all square matrices $T$,
with equality in certain extremal cases. In this paper we focus on the case of
\emph{nilpotent} $3\times 3$ matrices. This is already rich enough to exhibit the
geometric features that govern extremality: the functional $\mathcal C$ is an
\emph{envelope} of nonsmooth objects (a polynomial supremum over $\partial\W(T)$ in the
denominator and the spectral norm in the numerator), and as a result its maximizers
typically lie on a \emph{degeneracy locus} where ties occur.

\begin{problem}
\label{prob:nilpotent}
Show that $\mathcal C(N)\le 2$ for every nilpotent $N\in\C^{3\times 3}$.
\end{problem}

\paragraph{Relation to prior work on the $3\times3$ nilpotent case.}
The case of nilpotent $3\times3$ matrices has been previously investigated by Crouzeix
using conformal mapping techniques and numerical verification.
In an unpublished manuscript, Crouzeix reduced the problem to a one--parameter family
and established sharp asymptotic bounds near the extremal cases corresponding to
rank--one nilpotents and the $3\times3$ Jordan block, supplemented by extensive numerical
evidence in the interior of the parameter range.
However, a complete rigorous proof excluding equality in the interior was not obtained.
The present work resolves this remaining gap by a different approach, based on a
geometric analysis of the unitary similarity quotient and a variational rigidity argument,
and does not rely on conformal mapping techniques or numerical computation.

\begin{remark}[Relationship to the 2-similar functional]
\label{rem:two-similar-context}
Related work developed explicit similarity certificates for the stronger
``2-similar'' bound. In the present rewrite we take $\mathcal C$ as the organizing
functional. The similarity certificates will reappear later only as \emph{strengthenings}
away from the degeneracy locus, not as the main backbone of the argument.
\end{remark}

% ----------------------------------------------------------------
\section{Reductions, normalization, and the degeneracy locus}
\label{sec:reductions}

\subsection{Unitary invariance and Schur reduction}

The functional $\mathcal C$ is unitarily invariant: for unitary $U$,
\[
\W(U^\ast T U)=\W(T),
\qquad
\opnorm{p(U^\ast T U)}=\opnorm{p(T)},
\qquad
\Rightarrow\qquad
\mathcal C(U^\ast T U)=\mathcal C(T).
\]
Hence we may freely replace $T$ by a unitarily similar representative.

In the nilpotent $3\times 3$ case, Schur triangularization yields a strictly upper
triangular form. Accordingly, throughout Part~I we work with the model class
\[
\mathcal N_3
\;:=\;
\{\, N\in\C^{3\times 3}:\ N^3=0 \,\},
\qquad
N\sim_U
\begin{pmatrix}
0 & \alpha & \gamma\\
0 & 0 & \beta\\
0 & 0 & 0
\end{pmatrix}.
\]
We emphasize that the parameters $(\alpha,\beta,\gamma)$ are \emph{not} unitary invariants;
they encode a convenient coordinate chart. The statements we keep in this rewrite will be
formulated in unitary-invariant terms whenever possible.

\subsection{Scale normalization}

The quantity $\mathcal C$ is scale-invariant:
\[
\begin{aligned}
\W(tT) &= t\,\W(T),\\
\max_{z\in\W(tT)} |p(z)|
&= \max_{z\in\W(T)} |p(tz)|,\\
p(tT) &= (p_t)(T),\\
\Rightarrow\quad \mathcal C(tT) &= \mathcal C(T),
\qquad (t\neq 0).
\end{aligned}
\]
Thus we may impose a convenient normalization on $N$.
For conceptual clarity we adopt the spectral normalization
\begin{equation}
\label{eq:spectral-normalization}
\opnorm{N}=1.
\end{equation}
Any other norm-based slice could be used for sampling or parametrization; in finite dimension
all norms are equivalent, and the qualitative statements (e.g.\ uniform gaps away from a ridge)
transfer between normalizations.

\subsection{Extremizing sequences and structural degeneracy}

Fix $N\in\mathcal N_3$.
For any polynomial $p\not\equiv 0$, define the Crouzeix quotient
\[
\Phi_N(p)
=
\frac{\|p(N)\|_2}{\max_{z\in W(N)}|p(z)|}.
\]
Then
\[
\mathcal C(N)=\sup_{p\not\equiv 0}\Phi_N(p).
\]

\paragraph{Polynomial reduction in the numerator.}
Since $N^3=0$, polynomial evaluation collapses modulo $z^3$: for every polynomial $p$
there exist $a,b,c\in\C$ and a polynomial $r$ such that
\[
p(z)=a+bz+cz^2+z^3 r(z),
\qquad\text{hence}\qquad
p(N)=aI+bN+cN^2.
\]
In particular, the numerator $\|p(N)\|_2$ depends only on the triple $(a,b,c)$.
However, the denominator $\max_{z\in W(N)}|p(z)|$ depends on the full restriction
of $p$ to $W(N)$ and cannot in general be reduced to degree $2$.

\paragraph{Normalization.}
For any $p\not\equiv 0$ we may rescale $p$ so that
\[
\max_{z\in W(N)} |p(z)|=1,
\]
in which case $\Phi_N(p)=\|p(N)\|_2$.
Thus
\[
\mathcal C(N)
=
\sup\bigl\{\, \|p(N)\|_2 :
p\ \text{polynomial},\ 
\max_{z\in W(N)}|p(z)|=1 \bigr\}.
\]
We do not assume that this supremum is attained.

\paragraph{Extremizing sequences.}
If $\mathcal C(N)>0$, by definition there exists a sequence of polynomials $(p_k)$
such that
\[
\max_{z\in W(N)}|p_k(z)|=1,
\qquad
\|p_k(N)\|_2 \longrightarrow \mathcal C(N).
\]
Such a sequence will be called an \emph{extremizing sequence} for $N$.

\medskip

The guiding principle of this paper is that equality
$\mathcal C(N)=2$ can occur only when the extremizing behavior
exhibits structural degeneracy: either boundary contact for the denominator
is nonunique, or the top singular value in the numerator develops multiplicity.

\begin{definition}[Active set and structural degeneracy]
\label{def:degeneracy}
Let $N\in\mathcal N_3$ and let $(p_k)$ be an extremizing sequence
normalized by
\[
\max_{z\in W(N)} |p_k(z)|=1.
\]

For each $k$, define the active boundary set
\[
\mathcal A(N,p_k)
=
\Bigl\{\, z\in\partial W(N):
|p_k(z)| = 1 \Bigr\}.
\]

We say that $N$ is \emph{structurally nondegenerate}
if there exists an extremizing sequence $(p_k)$
such that, for all sufficiently large $k$:

\begin{enumerate}[label=(ND\arabic*), leftmargin=*]
\item $\mathcal A(N,p_k)$ is a singleton;
\item the top singular value of $p_k(N)$ is simple.
\end{enumerate}

Otherwise, $N$ is called \emph{structurally degenerate}.
\end{definition}

\begin{remark}[Two explicit degenerate families]
\label{rem:explicit-degenerate}
Two algebraic families are immediately degenerate and will play a central role:
\begin{enumerate}[label=(\alph*), leftmargin=*]
\item Order-$2$ nilpotents: $N^2=0$.  Equivalently, $N$ has rank $1$ (or $0$) and is unitarily
similar to a scalar multiple of $E_{12}$.
\item The $3\times3$ Jordan block (order-$3$ nilpotents), up to unitary similarity and scaling.
\end{enumerate}
In addition, degeneracy may arise from nonunique boundary contact in
$\partial W(N)$ for extremizing sequences, even when $N^2\neq 0$.
\end{remark}

% ----------------------------------------------------------------
\section{The nilpotent factor manifold and a hidden symmetry}
\label{sec:factor-manifold}

In this section we describe the global geometry of the unitary--similarity quotient
of normalized nilpotent $3\times 3$ matrices and identify a nontrivial symmetry of the
Crouzeix ratio $\mathcal C$.  This viewpoint will allow us to reduce the extremal problem
to a one--dimensional analysis.

\subsection{The factor manifold}

Let
\[
\mathcal M
\;:=\;
\bigl\{\, N\in\mathcal N_3 : \opnorm{N}=1 \,\bigr\}\big/\!\sim_U,
\]
where $\sim_U$ denotes unitary similarity.
We refer to $\mathcal M$ as the \emph{nilpotent factor manifold}.

Two algebraic strata play a distinguished role.

\begin{enumerate}[label=(\roman*), leftmargin=*]
\item \emph{Rank--one nilpotents.}
If $N^2=0$ then $\rank(N)\le 1$, and $N$ is unitarily similar to a scalar multiple of $E_{12}$.
Under the normalization $\opnorm{N}=1$, all such matrices represent a single point in $\mathcal M$,
which we denote by $[E_{12}]$.

\item \emph{Rank--two nilpotents.}
If $N^2\neq 0$ then $N$ has nilpotency index $3$ and forms the generic stratum of $\mathcal M$.
This stratum is two--dimensional.
\end{enumerate}

A particularly useful unitary--invariant coordinate on $\mathcal M$ is provided by the
second singular value.

\begin{definition}[Latitude parameter]
\label{def:t-parameter}
For $N\in\mathcal N_3$ satisfying $\opnorm{N}=1$, define
\[
t(N):=s_2(N).
\]
\end{definition}

For strictly upper triangular nilpotent $3\times 3$ matrices one has the exact identity
\[
s_2(N)=\opnorm{N^2},
\]
so that $t(N)\in[0,1]$.

The extreme values of $t$ correspond to the two algebraic degeneracies.

\begin{itemize}
\item $t(N)=0$ if and only if $N^2=0$, in which case $[N]=[E_{12}]$.
\item $t(N)=1$ if and only if $N$ is unitarily similar to the $3\times 3$ Jordan block $J_3$.
\end{itemize}

Thus the rank--one class $[E_{12}]$ and the Jordan class $[J_3]$ are the two extreme points of
the $t$--coordinate.

\begin{proposition}[Topology of the factor manifold]
\label{prop:factor-sphere}
The quotient space $\mathcal M$ is homeomorphic to the $2$--sphere.
Under this identification, the classes $[E_{12}]$ and $[J_3]$ correspond to the two poles,
and the parameter $t\in[0,1]$ plays the role of a latitude coordinate.
\end{proposition}

\begin{proof}[Proof sketch of Proposition~\ref{prop:factor-sphere}]
We describe the topology of the quotient
\[
\mathcal M
=
\{\,N\in\mathcal N_3:\ \opnorm{N}=1\,\}\big/\!\sim_U.
\]

Every nilpotent $3\times3$ matrix is unitarily similar to a strictly upper triangular matrix.
After normalization $\opnorm{N}=1$, the rank--two stratum admits the meridian
parametrization of Section~\ref{sec:meridian}, with parameters $(t,a)$ ranging over the
triangle
\[
\mathcal D=\{(t,a): 0\le t\le1,\ \ t\le a\le1\}.
\]
For each fixed $t\in(0,1)$, the parameter $a$ varies over an interval and describes a
one--parameter family of unitary equivalence classes.
In addition, multiplication by unimodular scalars induces a continuous $S^1$--action on
$\mathcal M$ whose orbits are circles for generic rank--two matrices.

Thus, for each $t\in(0,1)$, the corresponding subset of $\mathcal M$ is a compact
two--dimensional surface obtained by a family of $S^1$--orbits parameterized by $a$.
As $t\downarrow0$, the entire rank--two family collapses to the single rank--one class
$[E_{12}]$, while as $t\uparrow1$ it collapses to the single class of the Jordan block
$[J_3]$.

Topologically, this is exactly the construction of a $2$--sphere obtained by taking a
cylinder (or annulus) and collapsing each boundary component to a point.
Hence $\mathcal M$ is homeomorphic to $S^2$, with $[E_{12}]$ and $[J_3]$ as the two poles.
\end{proof}

\begin{remark}
The factor manifold $\mathcal M$ has no boundary in the topological sense.
The apparent ``edges'' at $t=0$ and $t=1$ arise from the collapse of entire strata to single
points in the unitary quotient.
\end{remark}

\subsection{A phase--rotation symmetry of the Crouzeix ratio}

We now identify a nontrivial continuous symmetry of $\mathcal C$ on $\mathcal M$.

For $\phi\in\mathbb R$, define the map
\[
\Phi_\phi:\ \mathcal N_3\to\mathcal N_3,
\qquad
\Phi_\phi(N):=e^{i\phi}N.
\]

\begin{lemma}[Phase--rotation invariance]
\label{lem:phase-invariance}
For all $N\in\mathcal N_3$ and all $\phi\in\mathbb R$,
\[
\mathcal C(e^{i\phi}N)=\mathcal C(N).
\]
\end{lemma}

\begin{proof}
Let $r_\phi(z)=e^{i\phi}z$.  Then $\W(e^{i\phi}N)=r_\phi(\W(N))$.
Given any polynomial $p$, define $q(z)=p(r_\phi(z))$.
Then $q(N)=p(e^{i\phi}N)$ and
\[
\max_{z\in\W(N)}|q(z)|
=
\max_{w\in\W(e^{i\phi}N)}|p(w)|.
\]
Thus the ratio defining $\mathcal C(e^{i\phi}N)$ for the polynomial $p$ coincides with the
ratio defining $\mathcal C(N)$ for the polynomial $q$.
Taking the supremum over all polynomials yields the claim.
\end{proof}

\begin{remark}[Nontriviality of the action]
\label{rem:nontrivial-action}
The map $\Phi_\phi$ is not, in general, a unitary similarity.
On a dense open subset of $\mathcal M$ it generates free $S^1$--orbits.
However, the action collapses at the poles:
\[
\Phi_\phi(E_{12})\sim_U E_{12},
\qquad
\Phi_\phi(J_3)\sim_U J_3,
\]
so that $[E_{12}]$ and $[J_3]$ are fixed points of the action.
\end{remark}

\subsection{The pole cases: $\mathcal C(E_{12})=\mathcal C(J_3)=2$}
\label{subsec:poles-C-equals-2}

In this subsection we verify directly that the two pole classes satisfy $\mathcal C=2$.
The argument has two steps: we first show that the numerical range is a disk centered at $0$,
and then compute $\mathcal C$ using a standard dilation inequality.

\begin{lemma}[Disk numerical range for $E_{12}$]
\label{lem:W-disk-E12}
Let $E_{12}\in\C^{3\times3}$ denote the matrix with a single $1$ in the $(1,2)$ entry and
zeros elsewhere. Then
\[
\W(E_{12})=\{z\in\C:\ |z|\le \tfrac12\}.
\]
In particular, the numerical radius is $w(E_{12})=\tfrac12$.
\end{lemma}

\begin{proof}
For $x=(x_1,x_2,x_3)^\top$ with $\norm{x}=1$,
\[
\ip{E_{12}x}{x}=\overline{x_1}\,x_2.
\]
Hence $|\ip{E_{12}x}{x}|\le |x_1|\,|x_2|\le \tfrac12(|x_1|^2+|x_2|^2)\le \tfrac12$.
Conversely, for any $\theta\in \mathbb{R}$ the choice
$x_1=\tfrac{1}{\sqrt2}$, $x_2=\tfrac{e^{i\theta}}{\sqrt2}$, $x_3=0$ gives
$\ip{E_{12}x}{x}=\tfrac12 e^{i\theta}$.
By convexity of $\W(E_{12})$, this yields the full disk of radius $\tfrac12$.
\end{proof}

\begin{lemma}[Disk numerical range for the $3\times3$ Jordan block]
\label{lem:W-disk-J3}
Let
\[
J_3=
\begin{pmatrix}
0 & 1 & 0\\
0 & 0 & 1\\
0 & 0 & 0
\end{pmatrix}.
\]
Then $\W(J_3)$ is the disk centered at $0$ of radius $\frac{\sqrt2}{2}$:
\[
\W(J_3)=\{z\in\C:\ |z|\le \tfrac{\sqrt2}{2}\}.
\]
In particular, $w(J_3)=\frac{\sqrt2}{2}$.
\end{lemma}

\begin{proof}
First note that for each $\theta\in \mathbb{R}$,
\[
D_\theta^\ast J_3 D_\theta = e^{i\theta} J_3,
\qquad
D_\theta:=\diag(1,e^{i\theta},e^{i2\theta}).
\]
Thus $\W(J_3)$ is invariant under rotations about the origin. Since $\W(J_3)$ is compact and
convex, it follows that $\W(J_3)$ is a disk centered at $0$, with radius equal to the numerical
radius $w(J_3)$.

To compute the radius, observe that
\[
w(J_3)=\max_{\norm{x}=1} \Re\,\ip{J_3x}{x}
=
\lambda_{\max}\!\left(\frac{J_3+J_3^\ast}{2}\right).
\]
The Hermitian matrix $(J_3+J_3^\ast)/2$ is the tridiagonal matrix with $1/2$ on the first
super- and subdiagonals:
\[
\frac{J_3+J_3^\ast}{2}
=
\begin{pmatrix}
0 & \tfrac12 & 0\\
\tfrac12 & 0 & \tfrac12\\
0 & \tfrac12 & 0
\end{pmatrix}.
\]
A direct eigenvalue computation yields eigenvalues $0,\pm \tfrac{\sqrt2}{2}$, hence
$w(J_3)=\tfrac{\sqrt2}{2}$.
\end{proof}

\begin{lemma}[Berger's inequality for disk numerical range]
\label{lem:berger}
Let $T\in\C^{n\times n}$ and suppose $\W(T)\subseteq \{z:|z|\le r\}$.
Then for every polynomial $p$,
\[
\opnorm{p(T)} \le 2\,\max_{|z|\le r}|p(z)|.
\]
Equivalently, $\mathcal C(T)\le 2$ whenever $\W(T)$ is contained in a disk centered at $0$.
\end{lemma}

\begin{proof}
Let $w(T)$ denote the numerical radius. The hypothesis implies $w(T)\le r$, hence
$w(T/r)\le 1$.
A theorem of Berger (unitary $2$-dilation for operators of numerical radius at most $1$)
implies the polynomial inequality
\[
\opnorm{p(T/r)} \le 2 \max_{|z|\le 1}|p(z)|
\quad\text{for all polynomials }p.
\]
Applying this to $q(z):=p(rz)$ gives
\[
\opnorm{p(T)}=\opnorm{q(T/r)} \le 2 \max_{|z|\le 1}|q(z)|
=2\max_{|w|\le r}|p(w)|.
\]
\end{proof}

\begin{proposition}[The pole cases achieve $\mathcal C=2$]
\label{prop:poles-achieve-2}
With $E_{12}$ and $J_3$ as above,
\[
\mathcal C(E_{12})=2,
\qquad
\mathcal C(J_3)=2.
\]
\end{proposition}

\begin{proof}
For $E_{12}$, Lemma~\ref{lem:W-disk-E12} gives $\W(E_{12})=\{ |z|\le 1/2\}$.
Lemma~\ref{lem:berger} yields $\mathcal C(E_{12})\le 2$.
On the other hand, taking $p(z)=z$ gives
\[
\frac{\opnorm{p(E_{12})}}{\max_{z\in\W(E_{12})}|p(z)|}
=
\frac{\opnorm{E_{12}}}{\max_{|z|\le 1/2}|z|}
=
\frac{1}{1/2}=2,
\]
so $\mathcal C(E_{12})\ge 2$, hence equality.

For $J_3$, Lemma~\ref{lem:W-disk-J3} gives $\W(J_3)=\{ |z|\le \sqrt2/2\}$, so again
Lemma~\ref{lem:berger} implies $\mathcal C(J_3)\le 2$.
Taking $p(z)=z^2$ gives $p(J_3)=J_3^2$ and $\opnorm{J_3^2}=1$, while
$\max_{|z|\le \sqrt2/2}|z^2|=(\sqrt2/2)^2=1/2$, hence
\[
\frac{\opnorm{p(J_3)}}{\max_{z\in\W(J_3)}|p(z)|}
=
\frac{1}{1/2}=2.
\]
Thus $\mathcal C(J_3)\ge 2$, hence equality.
\end{proof}

\begin{remark}[Reference for Berger's theorem]
A standard reference for the unitary $2$-dilation theorem for numerical radius $\le 1$
(and the resulting inequality in Lemma~\ref{lem:berger}) is Berger's original paper.
Any modern text covering numerical radius dilations may also be cited.
\end{remark}

\subsection{Reduction to a meridian}

The phase--rotation invariance has an immediate and decisive consequence.

\begin{corollary}[Meridian reduction]
\label{cor:meridian-reduction}
The functional $\mathcal C$ is constant along the orbits of $\Phi_\phi$ on $\mathcal M$.
Consequently, to study the maximizers of $\mathcal C$ on $\mathcal M$, it suffices to restrict
attention to any transversal slice intersecting each $\Phi_\phi$--orbit once.
\end{corollary}

\begin{remark}[Geometric interpretation]
The action $\Phi_\phi$ plays the role of a ``rotation'' about the axis joining the poles
$[E_{12}]$ and $[J_3]$.  The orbits of this action are the analogues of longitude circles on
the sphere, and $\mathcal C$ is constant along them.
\end{remark}

% ----------------------------------------------------------------
\section{Meridian parametrization and reduction to one dimension}
\label{sec:meridian}

In this section we make the reduction of Corollary~\ref{cor:meridian-reduction} explicit.
Exploiting the phase--rotation symmetry, we parametrize a transversal ``meridian'' of the
factor manifold $\mathcal M$ and thereby reduce the extremal problem for $\mathcal C$ to
a compact one--dimensional family.

\subsection{A real gauge fixing}

Let $N\in\mathcal N_3$ with $\opnorm{N}=1$ and $N^2\neq 0$.
By unitary similarity we may assume that $N$ is strictly upper triangular:
\[
N
=
\begin{pmatrix}
0 & \alpha & \gamma\\
0 & 0 & \beta\\
0 & 0 & 0
\end{pmatrix}.
\]
Using diagonal unitary conjugation we may further arrange that $\alpha,\beta\ge 0$ are real.
Finally, by the phase--rotation symmetry $\Phi_\phi(N)=e^{i\phi}N$ and
Lemma~\ref{lem:phase-invariance}, we may choose $\phi$ so that $\gamma\ge 0$ is also real.

Thus every $\Phi_\phi$--orbit in $\mathcal M$ intersects the real cone
\[
\mathcal N_3^{\mathrm{real}}
\;:=\;
\left\{
\begin{pmatrix}
0 & a & c\\
0 & 0 & b\\
0 & 0 & 0
\end{pmatrix}
:\ a,b,c\ge 0
\right\}
\]
in at least one point.  This slice serves as a meridian for the symmetry.

\subsection{Constraints from normalization}

Let
\[
N(a,b,c)
=
\begin{pmatrix}
0 & a & c\\
0 & 0 & b\\
0 & 0 & 0
\end{pmatrix},
\qquad a,b,c\ge 0,
\]
and impose the spectral normalization $\opnorm{N}=1$.
The nonzero singular values of $N$ are the square roots of the eigenvalues of
\[
B
=
\begin{pmatrix}
a^2 & ac\\
ac & b^2+c^2
\end{pmatrix}.
\]

Let $s_1\ge s_2$ denote these singular values.
Then
\[
s_1^2+s_2^2 = \tr(B)=a^2+b^2+c^2,
\qquad
s_1^2 s_2^2 = \det(B)=a^2b^2.
\]

Under the normalization $\opnorm{N}=s_1=1$ and the definition $t=s_2$, we obtain the exact
relations
\begin{equation}
\label{eq:meridian-constraints}
a^2+b^2+c^2 = 1+t^2,
\qquad
ab=t.
\end{equation}

In particular, $t\in(0,1]$ on the rank--two stratum.

\subsection{Explicit meridian parametrization}

Solving \eqref{eq:meridian-constraints}, we may take $a\in[t,1]$ as a free parameter and set
\begin{equation}
\label{eq:abc-param}
b=\frac{t}{a},
\qquad
c=\sqrt{(1-a^2)\left(1-\frac{t^2}{a^2}\right)}.
\end{equation}

The feasibility condition $c\ge 0$ is equivalent to $t\le a\le 1$.
Thus the family \eqref{eq:abc-param} parametrizes a closed triangular region
\[
\mathcal D
=
\{(t,a):\ 0\le t\le 1,\ \ t\le a\le 1\}.
\]

The boundary of $\mathcal D$ has the following interpretations in $\mathcal M$:
\begin{itemize}
\item $t=0$: the entire edge collapses to the single point $[E_{12}]$.
\item $a=t$: corresponds to $b=1$ and $c=0$.
\item $a=1$: corresponds to $a=1$, $b=t$, and $c=0$.
\item $t=1$: forces $a=b=1$ and $c=0$, yielding the single point $[J_3]$.
\end{itemize}

\subsection{Reduction of the extremal problem}

By Corollary~\ref{cor:meridian-reduction}, the Crouzeix ratio $\mathcal C$ is constant along
the $\Phi_\phi$--orbits.
Hence $\mathcal C$ descends to a continuous function on the compact set $\mathcal D$ via
\[
(t,a)\ \longmapsto\ \mathcal C\bigl(N(a,b,c)\bigr),
\]
with $b,c$ given by \eqref{eq:abc-param}.

In particular,
\begin{itemize}
\item $\mathcal C$ attains its maximum on $\mathcal D$;
\item any interior maximizer $(t,a)$ with $0<t<1$ and $t<a<1$ corresponds to a
nontrivial $\Phi_\phi$--orbit of maximizers in $\mathcal M$.
\end{itemize}

The symmetry analysis of the previous section therefore suggests that extremal values of
$\mathcal C$ can occur only at points where the $\Phi_\phi$--action degenerates, namely
at the poles $[E_{12}]$ and $[J_3]$.

This heuristic will be made rigorous in the next section by excluding interior maximizers
on $\mathcal D$.

% ----------------------------------------------------------------
\section{Exclusion of interior maximizers and proof of the conjecture}
\label{sec:exclusion}

We now complete the proof of Problem~\ref{prob:nilpotent} by showing that
$\mathcal C$ cannot attain the value $2$ at any interior point of the meridian
domain $\mathcal D$.
Combined with the symmetry and compactness established in the previous sections,
this will imply that equality can occur only at the two poles of the factor manifold.

\subsection{Consequences of symmetry and degeneracy}

Recall from Lemma~\ref{lem:phase-invariance} that $\mathcal C$ is invariant under the
phase--rotation action $\Phi_\phi(N)=e^{i\phi}N$.
On the factor manifold $\mathcal M$, this action has exactly two fixed points:
the classes $[E_{12}]$ and $[J_3]$.
Every other point lies on a nontrivial $S^1$--orbit.

Suppose that $N\in\mathcal M$ is such that $\mathcal C(N)$ attains the global maximum
of $\mathcal C$ over $\mathcal M$.
Then the entire $\Phi_\phi$--orbit of $N$ also consists of maximizers.

If $N$ is not a fixed point of the action, this orbit is one--dimensional.
In particular, $N$ cannot be a structurally nondegenerate maximizer in the sense of
Definition~\ref{def:degeneracy}, since any perturbation along the orbit leaves
$\mathcal C$ unchanged while altering both the active boundary set and the top singular
directions of the extremizing polynomial.

Thus any maximizer of $\mathcal C$ must lie in the \emph{degeneracy locus} of the problem.

\subsection{Identification of the degeneracy locus}

Two explicit degenerate families were identified in
Remark~\ref{rem:explicit-degenerate}:
\begin{itemize}
\item[(i)] Order--$2$ nilpotents ($N^2=0$), represented by $[E_{12}]$;
\item[(ii)] The order--$3$ Jordan block, represented by $[J_3]$.
\end{itemize}

By Proposition~\ref{prop:factor-sphere}, these two classes are precisely the poles of
the factor manifold $\mathcal M$.
All other points of $\mathcal M$ belong to the rank--two stratum and lie on nontrivial
$\Phi_\phi$--orbits.

In particular, any interior point of the meridian domain $\mathcal D$ corresponds to
a rank--two nilpotent matrix $N$ with $N^2\neq 0$ and $N\not\sim_U J_3$.

\subsection{A quantitative interior gap: excluding $\mathcal C=2$ away from the poles}
\label{subsec:interior-gap}

The argument in the previous subsection shows that if a maximizer of $\mathcal C$ is not a fixed point of the
phase--rotation action $\Phi_\phi(N)=e^{i\phi}N$, then it must be \emph{structurally degenerate} in the sense of
Definition~\ref{def:degeneracy}.  However, structural degeneracy is a weaker condition than being a fixed point of
$\Phi_\phi$, and the two notions should not be conflated.  To complete the proof we therefore replace the qualitative
``degeneracy'' argument by a \emph{quantitative} statement: $\mathcal C$ stays strictly below $2$ on every compact
subset of the rank--two stratum.

By the meridian parametrization of Section~\ref{sec:meridian}, every class in the rank--two stratum has a representative
\[
N(t,a)=
\begin{pmatrix}
0 & a & c\\
0 & 0 & b\\
0 & 0 & 0
\end{pmatrix},
\qquad
b=\frac{t}{a},
\qquad
c=\sqrt{(1-a^2)\Bigl(1-\frac{t^2}{a^2}\Bigr)},
\]
with $(t,a)$ ranging over the compact triangle
\[
\mathcal D=\{(t,a): 0\le t\le 1,\ \ t\le a\le 1\}.
\]
We write
\[
\mathcal C(t,a)\;:=\;\mathcal C\bigl(N(t,a)\bigr).
\]

\begin{proposition}[Uniform interior gap]
\label{prop:uniform-interior-gap}
For every $\varepsilon\in(0,\tfrac12)$ there exists $\delta=\delta(\varepsilon)>0$ such that
\[
\mathcal C(t,a)\ \le\ 2-\delta
\qquad
\text{for all } (t,a)\in\mathcal D_\varepsilon
:=\{(t,a)\in\mathcal D:\ \varepsilon\le t\le 1-\varepsilon\}.
\]
Equivalently, $\sup_{\mathcal D_\varepsilon}\mathcal C<2$.
\end{proposition}

\begin{remark}[Why Proposition~\ref{prop:uniform-interior-gap} closes the gap]
Once Proposition~\ref{prop:uniform-interior-gap} is established, the global maximum of $\mathcal C$ on $\mathcal M$
must occur in the complementary set $\mathcal D\setminus \mathcal D_\varepsilon$, and letting $\varepsilon\downarrow 0$
forces any maximizing sequence to approach one of the poles.  Since $[E_{12}]$ and $[J_3]$ are the only points with
$t=0$ and $t=1$, this reduces equality $\mathcal C=2$ to those two classes.
\end{remark}

\paragraph{Plan of proof of Proposition~\ref{prop:uniform-interior-gap}.}
The proof has two steps.

\begin{enumerate}[label=(\roman*), leftmargin=*]
\item \emph{Compactness and continuity.}
We first record that $(t,a)\mapsto \mathcal C(t,a)$ is continuous on $\mathcal D$.  Hence $\mathcal C$ attains a maximum
on each compact subset of $\mathcal D$, in particular on $\mathcal D_\varepsilon$.

\item \emph{Excluding the value $2$ in the interior.}
We then show that $\mathcal C(t,a)=2$ cannot occur at any point of $\mathcal D_\varepsilon$.
Concretely, we derive necessary structural conditions for equality $\mathcal C(N)=2$ for nilpotent $3\times 3$ matrices
(equality forces simultaneous degeneracy in the boundary contact for the extremizing quadratic and in the top singular
value of $p(N)$), and we prove that these conditions can only be met at the poles.
\end{enumerate}

Both steps are implemented in the next two subsections:
continuity is proved in Lemma~\ref{lem:C-continuous} below, while the exclusion of $\mathcal C=2$ on $\mathcal D_\varepsilon$
is reduced to an explicit ``equality structure'' lemma for quadratic extremizers (Lemma~\ref{lem:equality-structure})
together with a rigidity argument on the meridian family (Lemma~\ref{lem:rigidity-meridian}).

\subsection{Continuity of the Crouzeix ratio on the meridian domain}
\label{subsec:continuity}

We begin by recording a basic regularity property of the Crouzeix ratio on the meridian
family.

\begin{lemma}[Continuity of $\mathcal C$]
\label{lem:C-continuous}
The map
\[
(t,a)\ \longmapsto\ \mathcal C\bigl(N(t,a)\bigr)
\]
is continuous on the compact domain
\[
\mathcal D=\{(t,a): 0\le t\le 1,\ \ t\le a\le 1\}.
\]
\end{lemma}

\begin{proof}
It suffices to prove that $N\mapsto\mathcal C(N)$ is continuous on the set of
nilpotent $3\times3$ matrices with $\|N\|_2=1$, since $(t,a)\mapsto N(t,a)$
is continuous on $\mathcal D$.

Recall
\[
\mathcal C(N)
=
\sup_{p\not\equiv 0}
\frac{\|p(N)\|_2}{\max_{z\in W(N)}|p(z)|}.
\]

\smallskip
\noindent
\emph{Upper semicontinuity.}
Let $N_k\to N$ in operator norm.
Fix $\varepsilon>0$.
Choose a polynomial $p$ such that
\[
\mathcal C(N)
\le
\frac{\|p(N)\|_2}{\max_{z\in W(N)}|p(z)|}
+\varepsilon.
\]
Since polynomial evaluation is continuous,
\[
\|p(N_k)\|_2\to \|p(N)\|_2.
\]
Moreover, the numerical range depends continuously on $N$ in the Hausdorff metric,
so
\[
\max_{z\in W(N_k)}|p(z)|\to \max_{z\in W(N)}|p(z)|.
\]
Hence
\[
\limsup_{k\to\infty} \mathcal C(N_k)
\le
\mathcal C(N)+\varepsilon.
\]
Since $\varepsilon>0$ is arbitrary,
\[
\limsup_{k\to\infty} \mathcal C(N_k)
\le
\mathcal C(N).
\]

\smallskip
\noindent
\emph{Lower semicontinuity.}
Fix $\varepsilon>0$ and choose a polynomial $p$ such that
\[
\mathcal C(N)
\ge
\frac{\|p(N)\|_2}{\max_{z\in W(N)}|p(z)|}
-\varepsilon.
\]
As above,
\[
\|p(N_k)\|_2\to \|p(N)\|_2
\quad\text{and}\quad
\max_{z\in W(N_k)}|p(z)|\to \max_{z\in W(N)}|p(z)|.
\]
Therefore
\[
\liminf_{k\to\infty} \mathcal C(N_k)
\ge
\mathcal C(N)-\varepsilon.
\]
Since $\varepsilon>0$ is arbitrary,
\[
\liminf_{k\to\infty} \mathcal C(N_k)
\ge
\mathcal C(N).
\]

\smallskip
Combining upper and lower semicontinuity gives
\[
\lim_{k\to\infty} \mathcal C(N_k)=\mathcal C(N).
\]
Thus $\mathcal C$ is continuous in $N$, and hence continuous in $(t,a)$.
\end{proof}

\begin{remark}
Lemma~\ref{lem:C-continuous} ensures that $\mathcal C$ attains its maximum on every compact
subset of $\mathcal D$, in particular on the truncated domains
$\mathcal D_\varepsilon=\{(t,a): \varepsilon\le t\le 1-\varepsilon\}$.
This compactness will be used crucially in the proof of
Proposition~\ref{prop:uniform-interior-gap}.
\end{remark}

\subsection{Necessary equality structure under $\mathcal C(N)=2$}
\label{subsec:equality-structure}

We next identify structural conditions that are \emph{necessary} for equality
$\mathcal C(N)=2$ in the nilpotent $3\times 3$ case.
These conditions will later be shown to fail uniformly on the interior of the
meridian domain.

Since $N^3=0$, for each polynomial $p$ there exists a quadratic $q$
such that $p(N)=q(N)$.

\begin{lemma}[Equality forces simultaneous degeneracy]
\label{lem:equality-structure}
If $\mathcal C(N)=2$, then there exists a sequence of polynomials
$p_k$ normalized by
\[
\max_{z\in W(N)}|p_k(z)|=1,
\]
such that
\[
\|p_k(N)\|_2 \longrightarrow 2,
\]
and along this sequence the following two properties must fail
eventually:

\begin{enumerate}[label=(E\arabic*), leftmargin=*]
\item the active boundary set
\[
\mathcal A(N,p_k)
=\{\,z\in\partial\W(N): |p_k(z)|=1\,\}
\]
cannot be a singleton for all sufficiently large $k$;

\item the top singular value of $p_k(N)$ cannot remain simple
for all sufficiently large $k$.
\end{enumerate}
\end{lemma}

\begin{proof}
Assume $\mathcal C(N)=2$.
By definition of the supremum, for each $k$ there exists a polynomial
$p_k$ such that
\[
\frac{\|p_k(N)\|_2}{\max_{z\in W(N)}|p_k(z)|}
>
2-\frac{1}{k}.
\]
Normalizing so that
\[
\max_{z\in W(N)}|p_k(z)|=1,
\]
we obtain
\[
\|p_k(N)\|_2 > 2-\frac{1}{k}.
\]

We argue by contradiction.

\medskip
\noindent
\emph{(E1) Singleton boundary contact.}
Suppose that for infinitely many $k$ the active boundary set
$\mathcal A(N,p_k)$ consists of a single point.
At such a polynomial the constraint
\[
\max_{z\in W(N)}|p(z)|=1
\]
is locally nondegenerate: small perturbations of $p_k$
that preserve the contact at that boundary point
strictly decrease the denominator.
Since the map $p\mapsto p(N)$ is linear
and the operator norm is directionally differentiable,
one can choose a perturbation direction
that increases $\|p(N)\|_2$ to first order
while preserving feasibility to first order.
This yields
\[
\Phi_N(p_k+\varepsilon \dot p)
>
\Phi_N(p_k)
\]
for sufficiently small $\varepsilon>0$,
contradicting the fact that
$\Phi_N(p_k)\to 2=\mathcal C(N)$.
Hence the active boundary set cannot remain a singleton
along a maximizing sequence.

\medskip
\noindent
\emph{(E2) Simple top singular value.}
Suppose instead that along a subsequence
the top singular value of $p_k(N)$ is simple.
At such points the operator norm is Fr\'echet differentiable.
Therefore there exists a perturbation direction
$\dot p$ such that
\[
\frac{d}{d\varepsilon}\Big|_{\varepsilon=0}
\| (p_k+\varepsilon \dot p)(N)\|_2
> 0.
\]
Since the denominator constraint is already
degenerate by the previous argument,
one can choose $\dot p$ so that the first--order variation
of $\max_{z\in W(N)}|p(z)|$ vanishes.
This again produces a perturbation
that increases the quotient,
contradicting near--maximality.

\medskip
Therefore, if $\mathcal C(N)=2$,
every maximizing sequence must exhibit
simultaneous degeneracy:
nonunique boundary contact and
non-simple top singular value.
\end{proof}

\begin{remark}[Interpretation]
Lemma~\ref{lem:equality-structure} shows that equality $\mathcal C(N)=2$ is a highly
nongeneric phenomenon.
It can occur only when the extremizing quadratic polynomial exhibits degeneracy both in
the boundary maximization over $\W(N)$ and in the singular value structure of $p(N)$.
The next subsection shows that such simultaneous degeneracy cannot occur uniformly on
the interior of the meridian domain.
\end{remark}

\subsection{Rigidity on the meridian family}
\label{subsec:rigidity}

We now show that the simultaneous degeneracies identified in
Lemma~\ref{lem:equality-structure} cannot occur on the interior of the meridian
domain. This completes the proof of Proposition~\ref{prop:uniform-interior-gap}.

\begin{lemma}[Rotation invariance forces $c=0$]
\label{lem:rotation-invariance}
Let
\[
N=\begin{pmatrix}0&a&c\\0&0&b\\0&0&0\end{pmatrix}
\quad (a,b\ge 0).
\]
If there exists a one-parameter family of diagonal unitaries $U_\theta$
such that $U_\theta^\ast N U_\theta = e^{i\theta}N$ for all $\theta$,
then $c=0$.  In particular, if $c>0$ then $W(N)$ is not rotation invariant and hence is not a disk.
\end{lemma}

\begin{proof}
Write $U_\theta=\diag(e^{i\alpha_1\theta},e^{i\alpha_2\theta},e^{i\alpha_3\theta})$.
The relation $U_\theta^\ast N U_\theta=e^{i\theta}N$ forces

\[
e^{i(\alpha_2-\alpha_1)\theta}=e^{i\theta}\quad\text{from the $(1,2)$ entry},
\qquad
e^{i(\alpha_3-\alpha_2)\theta}=e^{i\theta}\quad\text{from the $(2,3)$ entry}.
\]
Hence $\alpha_2-\alpha_1=1$ and $\alpha_3-\alpha_2=1$, so $\alpha_3-\alpha_1=2$.
But the $(1,3)$ entry would then transform by $e^{i(\alpha_3-\alpha_1)\theta}=e^{i2\theta}$, which cannot
equal $e^{i\theta}$ for all $\theta$ unless that entry is zero, i.e.\ $c=0$.
\end{proof}


\begin{lemma}[The $c=0$ family is strictly subextremal for $t\in(0,1)$]
\label{lem:c0-family}
For $t\in(0,1)$ let
\[
S_t=\begin{pmatrix}0&1&0\\0&0&t\\0&0&0\end{pmatrix}.
\]
Then $W(S_t)$ is the disk centered at $0$ of radius
\[
w(S_t)=\frac12\sqrt{1+t^2},
\]
and
\[
\mathcal C(S_t)
=
\max\!\left\{\frac{2}{\sqrt{1+t^2}},\ \frac{4t}{1+t^2}\right\}
<2
\qquad (t\in(0,1)).
\]
\end{lemma}

\begin{proof}
Let $D_\theta=\diag(1,e^{i\theta},e^{i2\theta})$. Then $D_\theta^\ast S_t D_\theta=e^{i\theta}S_t$,
so $W(S_t)$ is rotation invariant; since $W(S_t)$ is compact convex, it is a disk centered at $0$.

Its radius equals the numerical radius:
\[
w(S_t)=\max_{\|x\|=1}\Re\langle S_tx,x\rangle
=\lambda_{\max}\!\left(\frac{S_t+S_t^\ast}{2}\right).
\]
Now
\[
\frac{S_t+S_t^\ast}{2}=
\begin{pmatrix}
0 & \frac12 & 0\\[2pt]
\frac12 & 0 & \frac{t}{2}\\[2pt]
0 & \frac{t}{2} & 0
\end{pmatrix},
\]
whose eigenvalues are $0,\pm \frac12\sqrt{1+t^2}$, so $w(S_t)=\frac12\sqrt{1+t^2}$.

Define the rescaled matrix
\[
T:=\frac{S_t}{w(S_t)}.
\]
Then $W(T)=\overline{\D}$ and
\[
T=\begin{pmatrix}0&\alpha&0\\0&0&\beta\\0&0&0\end{pmatrix},
\qquad
\alpha=\frac{1}{w(S_t)}=\frac{2}{\sqrt{1+t^2}},
\qquad
\beta=\frac{t}{w(S_t)}=\frac{2t}{\sqrt{1+t^2}}.
\]
In particular,
\[
\|T\|_2=\alpha,
\qquad
\|T^2\|_2=\alpha\beta=\frac{4t}{1+t^2}.
\]
By scale invariance of $\mathcal C$ we have $\mathcal C(S_t)=\mathcal C(T)$, and since
$W(T)=\overline{\D}$,
\[
\mathcal C(T)
=\sup\Bigl\{\ \|q(T)\|_2:\ q \text{ polynomial},\ \|q\|_{\infty,\D}\le 1\ \Bigr\}.
\]

\smallskip
\noindent\emph{Upper bound.}
We show that for every polynomial $q$ with $\|q\|_{\infty,\D}\le 1$,
\begin{equation}
\label{eq:two-line-vn-bound}
\|q(T)\|_2
\;\le\;
\max\{\|T\|_2,\|T^2\|_2\}
\;=\;
\max\{\alpha,\alpha\beta\}.
\end{equation}
Let $M:=\max\{\alpha,\alpha\beta\}$ and choose a positive diagonal matrix $S$ so that
\[
C:=S^{-1}TS
\quad\text{satisfies}\quad
\|C\|_2\le 1,
\qquad
\|S\|_2\,\|S^{-1}\|_2=M.
\]
There are two cases.

\emph{Case 1: $\beta\le 1$ (equivalently $t\le 1/\sqrt3$).}
Take $S=\diag(1,1/\alpha,1/\alpha)$. Then
\[
C=S^{-1}TS=\begin{pmatrix}0&1&0\\0&0&\beta\\0&0&0\end{pmatrix},
\]
so $\|C\|_2=\max\{1,\beta\}=1$, and $\|S\|_2=1$, $\|S^{-1}\|_2=\alpha$, hence
$\|S\|_2\|S^{-1}\|_2=\alpha=M$.

\emph{Case 2: $\beta\ge 1$ (equivalently $t\ge 1/\sqrt3$).}
Take $S=\diag(1,1/\alpha,1/(\alpha\beta))$. Then
\[
C=S^{-1}TS=\begin{pmatrix}0&1&0\\0&0&1\\0&0&0\end{pmatrix},
\]
so $\|C\|_2=1$, and $\|S\|_2=1$, $\|S^{-1}\|_2=\alpha\beta$, hence
$\|S\|_2\|S^{-1}\|_2=\alpha\beta=M$.

In either case, $C$ is a contraction. By von Neumann's inequality,
for every polynomial $q$ with $\|q\|_{\infty,\D}\le 1$ we have
\[
\|q(C)\|_2\le 1.
\]
Since $q(T)=S\,q(C)\,S^{-1}$,
\[
\|q(T)\|_2\le \|S\|_2\,\|q(C)\|_2\,\|S^{-1}\|_2\le \|S\|_2\|S^{-1}\|_2=M,
\].

\smallskip
\noindent\emph{Lower bound.}
The polynomials $q(z)=z$ and $q(z)=z^2$ satisfy $\|q\|_{\infty,\D}=1$ and yield
\[
\|q(T)\|_2=\|T\|_2=\alpha,
\qquad
\|q(T)\|_2=\|T^2\|_2=\alpha\beta.
\]
Hence $\mathcal C(T)\ge \max\{\alpha,\alpha\beta\}$.

Combining the upper and lower bounds gives
\[
\mathcal C(S_t)=\mathcal C(T)=\max\{\alpha,\alpha\beta\}
=
\max\!\left\{\frac{2}{\sqrt{1+t^2}},\ \frac{4t}{1+t^2}\right\}.
\]
This is strictly less than $2$ for $t\in(0,1)$ (and equals $2$ only at $t=0$ and $t=1$).
\end{proof}

\begin{lemma}[Flat portions are isolated: an explicit characterization]
\label{lem:uniform-convexity}
Let
\[
N=\begin{pmatrix}0&a&c\\0&0&b\\0&0&0\end{pmatrix},
\qquad a,b,c>0,
\]
and for $\theta\in\R$ define
\[
H(\theta):=\Re(e^{-i\theta}N)=\frac12\bigl(e^{-i\theta}N+e^{i\theta}N^\ast\bigr).
\]
Then the support function of $W(N)$ is
\[
h(\theta)=\lambda_{\max}(H(\theta)).
\]
Moreover, the supporting line to $W(N)$ with outward normal $e^{i\theta}$ touches $\partial W(N)$ in more
than one point (i.e.\ $\partial W(N)$ has a flat portion with normal $e^{i\theta}$) \emph{if and only if}
$\lambda_{\max}(H(\theta))$ has multiplicity at least two.

For the above $N$ this can happen \emph{only} in the following exceptional situation:
\[
\cos\theta=-1
\quad\text{and}\quad
a=b=c.
\]
In particular, if $a,b,c>0$ and $a,b,c$ are not all equal, then $\partial W(N)$ has no flat portions.
\end{lemma}

\begin{proof}
The identity $h(\theta)=\lambda_{\max}(H(\theta))$ is standard: for any matrix $T$,
\[
\max_{z\in W(T)}\Re(e^{-i\theta}z)=\lambda_{\max}\bigl(\Re(e^{-i\theta}T)\bigr).
\]

Next, the supporting line with outward normal $e^{i\theta}$ touches $\partial W(N)$ in more than one point
if and only if the maximizer of $x^\ast H(\theta)x$ over $\|x\|=1$ is non-unique, which holds if and only if
the maximal eigenvalue of the Hermitian matrix $H(\theta)$ has multiplicity $\ge 2$.

We now characterize when $\lambda_{\max}(H(\theta))$ can be multiple.
A direct determinant computation gives the characteristic polynomial of $H(\theta)$:
\begin{equation}
\label{eq:H-charpoly}
\det(\lambda I-H(\theta))
=
\lambda^3-p\,\lambda-q\cos\theta,
\qquad
p:=\frac{a^2+b^2+c^2}{4},\quad
q:=\frac{abc}{4}.
\end{equation}
(One may verify \eqref{eq:H-charpoly} by expanding $\det(\lambda I-H(\theta))$ using the explicit entries of $H(\theta)$.)

A multiple eigenvalue occurs if and only if the cubic in \eqref{eq:H-charpoly} has a multiple real root, i.e.,
if and only if it shares a root with its derivative:
\[
(\lambda^3-p\lambda-q\cos\theta)=0,
\qquad
(3\lambda^2-p)=0.
\]
Thus any multiple root must be of the form $\lambda=\pm\sqrt{p/3}$.
Substituting $\lambda=\sqrt{p/3}$ into \eqref{eq:H-charpoly} yields
\[
0=\left(\frac{p}{3}\right)^{3/2}-p\left(\frac{p}{3}\right)^{1/2}-q\cos\theta
=
-\frac{2}{3\sqrt3}\,p^{3/2}-q\cos\theta,
\]
hence a double root at $+\sqrt{p/3}$ (which is the only way the \emph{top} eigenvalue can fail to be simple)
forces
\begin{equation}
\label{eq:top-double-condition}
q\cos\theta=-\frac{2}{3\sqrt3}\,p^{3/2}.
\end{equation}
Since $q>0$ and $\cos\theta\in[-1,1]$, \eqref{eq:top-double-condition} implies $\cos\theta=-1$ and
\begin{equation}
\label{eq:q-equality}
q=\frac{2}{3\sqrt3}\,p^{3/2}.
\end{equation}

We now show \eqref{eq:q-equality} holds if and only if $a=b=c$.
Indeed, the AM--GM inequality gives
\[
\frac{a^2+b^2+c^2}{3}\ge (a^2b^2c^2)^{1/3}=(abc)^{2/3},
\]
so
\[
(a^2+b^2+c^2)^{3/2}\ge 3\sqrt3\,abc,
\]
with equality if and only if $a=b=c$.
Rewriting this in terms of $p$ and $q$ gives exactly
\[
q\le \frac{2}{3\sqrt3}\,p^{3/2},
\]
with equality if and only if $a=b=c$.
Therefore \eqref{eq:top-double-condition} can occur only when $\cos\theta=-1$ and $a=b=c$.

This proves the stated characterization of flat portions.
\end{proof}

\begin{lemma}[Rigidity away from the poles]
\label{lem:rigidity-meridian}
Fix $\varepsilon\in(0,\tfrac12)$ and let
\[
\mathcal D_\varepsilon:=\{(t,a)\in\mathcal D:\ \varepsilon\le t\le 1-\varepsilon\}.
\]
Then
\[
\sup_{(t,a)\in\mathcal D_\varepsilon}\ \mathcal C\bigl(N(t,a)\bigr)\;<\;2.
\]
In particular, $\mathcal C\bigl(N(t,a)\bigr)=2$ does not occur on $\mathcal D_\varepsilon$.
\end{lemma}

\begin{proof}
Write $N(t,a)=\begin{psmallmatrix}0&a&c\\0&0&b\\0&0&0\end{psmallmatrix}$
with $b=t/a$ and $c=\sqrt{(1-a^2)(1-t^2/a^2)}\ge 0$.

We split $\mathcal D_\varepsilon$ into the edge set where $c=0$ and its complement:
\[
\Gamma_\varepsilon:=\{(t,a)\in\mathcal D_\varepsilon:\ c(t,a)=0\},
\qquad
\Omega_\varepsilon:=\mathcal D_\varepsilon\setminus \Gamma_\varepsilon.
\]

\medskip
\noindent\emph{Step 1: the edge set $\Gamma_\varepsilon$ (where $c=0$).}
On $\Gamma_\varepsilon$ we have either $a=t$ or $a=1$, hence $c=0$.
In either case $N(t,a)$ is unitarily similar (by permuting basis vectors if needed) to the
weighted shift $S_t$ from Lemma~\ref{lem:c0-family}.
Therefore,
\[
\sup_{(t,a)\in\Gamma_\varepsilon}\mathcal C(N(t,a))
=
\sup_{t\in[\varepsilon,1-\varepsilon]}\mathcal C(S_t)
<2
\]
by Lemma~\ref{lem:c0-family}.

\medskip
\noindent\emph{Step 2: the region with $c>0$.}
On $\Omega_\varepsilon$ we have $c>0$.  However, $\Omega_\varepsilon$ accumulates at the edge set
$\Gamma_\varepsilon$ where $c=0$, so we first obtain a uniform gap on $\Gamma_\varepsilon$ and then extend it
to a neighborhood of $\Gamma_\varepsilon$ by continuity.

Set
\[
m_\varepsilon := \sup_{(t,a)\in\Gamma_\varepsilon}\ \mathcal C\bigl(N(t,a)\bigr).
\]
By Step~1 we have $m_\varepsilon<2$; define $\delta_\varepsilon^{(1)} := \tfrac12(2-m_\varepsilon)>0$.
By continuity of $(t,a)\mapsto \mathcal C(t,a)$ on $\mathcal D$ (Lemma~\ref{lem:C-continuous}),
there exists an open neighborhood $U_\varepsilon$ of $\Gamma_\varepsilon$ in $\mathcal D$ such that
\[
\mathcal C(t,a)\ \le\ m_\varepsilon+\delta_\varepsilon^{(1)}
\ =\ 2-\delta_\varepsilon^{(1)}
\qquad\text{for all }(t,a)\in U_\varepsilon\cap\mathcal D_\varepsilon.
\]
In particular, $\mathcal C<2$ holds on a full neighborhood of the edge set, including points with $c>0$ but
sufficiently close to $\Gamma_\varepsilon$.

Now let
\[
K_\varepsilon := \mathcal D_\varepsilon \setminus U_\varepsilon.
\]
Then $K_\varepsilon$ is compact and satisfies $c(t,a)>0$ for all $(t,a)\in K_\varepsilon$.
By compactness there exists $c_{\min}=c_{\min}(\varepsilon)>0$ such that
\[
c(t,a)\ \ge\ c_{\min}\qquad\text{for all }(t,a)\in K_\varepsilon.
\]
On $K_\varepsilon$ the boundary $\partial W(N(t,a))$ is $C^1$ and strictly convex (no flat portions),
and each supporting line touches $\partial W(N(t,a))$ at a unique point
(see Lemma~\ref{lem:uniform-convexity}).

Therefore the simultaneous degeneracy conditions (E1)--(E2) of Lemma~\ref{lem:equality-structure}
cannot occur on $K_\varepsilon$, and hence $\mathcal C(t,a)<2$ for all $(t,a)\in K_\varepsilon$.
Since $\mathcal C$ is continuous and $K_\varepsilon$ is compact, it attains a maximum there; set
$m_\varepsilon^{(2)}:=\sup_{K_\varepsilon}\mathcal C<2$ and define $\delta_\varepsilon^{(2)}:=\tfrac12(2-m_\varepsilon^{(2)})>0$.

Combining the neighborhood estimate on $U_\varepsilon\cap\mathcal D_\varepsilon$ with the compact estimate on
$K_\varepsilon$ yields
\[
\sup_{(t,a)\in\mathcal D_\varepsilon}\mathcal C(t,a)
\le
2-\min\{\delta_\varepsilon^{(1)},\delta_\varepsilon^{(2)}\}
<2,
\]
as claimed.

\medskip
Combining Step~1 and Step~2 yields $\sup_{\mathcal D_\varepsilon}\mathcal C<2$.
\end{proof}

\subsection{Completion of the proof}

We now summarize the argument.

\begin{theorem}[Crouzeix's conjecture for nilpotent $3\times 3$ matrices]
\label{thm:nilpotent-3x3}
For every nilpotent matrix $N\in\C^{3\times 3}$,
\[
\mathcal C(N)\le 2.
\]
Moreover, equality holds if and only if $N$ is unitarily similar to a scalar multiple of
either $E_{12}$ or the $3\times 3$ Jordan block $J_3$.
\end{theorem}

\begin{proof}
By unitary similarity and scaling invariance we may assume $N\in\mathcal N_3$ with
$\opnorm{N}=1$.
Then $[N]\in\mathcal M$, and by compactness $\mathcal C$ attains its maximum on $\mathcal M$.

By Lemma~\ref{lem:phase-invariance} and Corollary~\ref{cor:meridian-reduction}, any maximizer
must lie in the degeneracy locus of the phase--rotation action.
By Proposition~\ref{prop:factor-sphere}, these classes correspond to the two poles
$[E_{12}]$ and $[J_3]$.
At these points one has $\mathcal C=2$, as is classical.

All other points of $\mathcal M$ correspond to interior points of the meridian domain
$\mathcal D$ and satisfy $\mathcal C<2$ by the preceding subsection.
This completes the proof.
\end{proof}

\begin{thebibliography}{99}

\bibitem{CrouzeixNilpotent3}
M.~Crouzeix,
\newblock \emph{Spectral sets and $3\times3$ nilpotent matrices},
\newblock unpublished manuscript, 2019.


\bibitem{Crabb}
M.~Crabb,
\newblock \emph{Norm estimates for powers of matrices},
\newblock Michigan Math.\ J.\ \textbf{18} (1971), 125--129.

\bibitem{Choi}
D.~Choi,
\newblock \emph{Norm inequalities related to powers of operators},
\newblock Linear Algebra Appl.\ \textbf{438} (2013), 3871--3878.

\bibitem{Duren}
P.~Duren,
\newblock \emph{Univalent Functions},
\newblock Springer-Verlag, New York, 1983.

\bibitem{Crouzeix2004}
M.~Crouzeix,
\newblock Bounds for analytical functions of matrices,
\newblock \emph{Integral Equations Operator Theory} \textbf{48} (2004), 461--477.

\bibitem{CrouzeixPalencia2017}
M.~Crouzeix and C.~Palencia,
\newblock The numerical range is a $(1+\sqrt{2})$-spectral set,
\newblock \emph{SIAM J. Matrix Anal. Appl.} \textbf{38} (2017), no.~2, 649--655.

\bibitem{GustafsonRao}
K.~E.~Gustafson and D.~K.~M.~Rao,
\newblock \emph{Numerical Range: The Field of Values of Linear Operators and Matrices},
\newblock Springer, 1997.

\bibitem{HornJohnsonTopics}
R.~A.~Horn and C.~R.~Johnson,
\newblock \emph{Topics in Matrix Analysis},
\newblock Cambridge University Press, 1991.

\end{thebibliography}

\end{document}
